% !TeX root = ../colloquium-2.tex

\phantomsection\label{sec_Proofs06}
\begin{proof}[Доказательство теоремы о производной композиции. Доказательство теоремы о производной обратной функции. Вычисление производных функции $f(x) = \arcsin x, f(x) = \ln x$]
    \subsection*{Доказательство теоремы о производной композиции.}
    \begin{note}[amber]
        Пусть функция $f$ дифференцируема в точке а, а функция $g$ дефференцируема в точке $f(a)$. Тогда функция $g \circ f$ дифференцируема в точке а и $(g \circ f)'(a) = g'(f(a)) \cdot f'(a)$\\
        \begin{note}[purple]
           Доказательства. Так как функция $g$ дифференцируемав точке $f(a)$, то при всех q из некоторой проколотой окрестности нуля
           \[
           g(f(a) + q) - g(f(a)) = g'(f(a))q + \beta (q)q, \lim_{q \to 0} \beta(q) = 0.
           \]
           Доопределим бесконечно малую функцию $\beta$ в нуле, положив $\beta(0) := 0$, то есть доопределим функцию $\beta$
           в нуле по непрерывности. Теперь мы можем положить $q = f(a+h) - f((a))$. При этом при подстановки вместо q, принадлежащего проколотой окрестности нуля , приращения $f(a+h)-f(a)$ может возникнуть ситуация, при которой приращение для $f$ ранво нулю.
           Для того, что бы при этом равенство для приращения $g$ не нарушалась, мы и доопределили бесконечно малую функцию $\beta$ в нуле по непрерывности. Отметим, что такое доопределение не нарушает дефференцируемость функции $g$,
           но позволяет корректно подставить $g = f(a+h)-f(a)$.
           Тогда в силу дифференцируемости функции $f$ из которой вытекает что:
           \[
           f(a+h)-f(a) = f'(a)h + \alpha(h)h, \lim_{h \to 0} \alpha(h) = 0.
           \]
           в точке $a$ получим:
           \[
            g(f(a+h))-g(f(a)) = 
            \]
            \[
            g'(f(a))(f(a+h)-f(a)) + \beta(f(a+h)-f(a))(f(a+h)-f(a)) =
            \]
            \[
            g'(f(a))f'(a)h+g'(f(a))\alpha(h)h+\beta(f'(a)h+\alpha(h)h)(f'(a)h+\alpha(h)h) =
            \]
            \[
            g'(f(a))f'(a)h + \delta (h)h.
           \]
           мы использовали обозначение $\delta(h) = g'(f(a))\alpha(h) + \beta(f'(a)h + \alpha(h)h)(f'(a) + \alpha(h))$,
           а тогда из арифметики пределов и из бесконечной малости функции $\beta$ и $\alpha$ следует, что $\delta(h) \rightarrow 0, h \rightarrow 0$. Таким образом, для функции $g \circ f$ в точке $a$ по определению. 
           Мы знаем, что это равносильно наличию у функции $g \circ f$ производной в точке $a$. Таким образом, получаем, что
           \[
            (g \circ f)'(a) = g'(f(a))f'(a).
           \]
        \end{note}
    \end{note}

    \subsection*{Доказательство теоремы о производной обратной функции.}
    \begin{note}[amber]
        Пусть функция
        
        \begin{note}[purple]
           
        \end{note}
    \end{note}
\end{proof}
